% !PRINTED 257
\weberpage{257}
%<<CONT
domain of rationality belong the functions of $\omega$ which, under
the two substitutions

\weq{1}{\left(\omega, -\frac{1}{\omega}\right), \quad (\omega, \omega+2)}

remain unchanged (\S\,54, 2). If, then, a system of $\nu$
functions $\Phi_{a,c,\partial}$ is permuted by the substitutions (1) only among
themselves, these functions are the roots of a
transformation equation whose coefficients depend rationally on
$f(\omega)^{24}$. To these belong (as we have already shown earlier by another
route), for an \gesperrt{odd} $n$, which we shall
now always assume, certain powers of the quantities

\[
f\left(\frac{c+\partial\omega}{a}\right),
\]

in which, as may be remarked once and for all, $c$ is assumed divisible by 16, and,
when $n$ is not divisible by 3, divisible by 48.

The somewhat extended principles of \S\,70 lead comparatively
simply to the calculation of these equations.

An extension is, however, necessary for the following reason:

In the considerations hitherto we were able to draw the
conclusion: if a \gesperrt{rational} function of $j(\omega)$ remains finite for every
finite $\omega$ with positive imaginary part, then it is
a \gesperrt{polynomial} in $j(\omega)$, because to every \gesperrt{finite} $j(\omega)$
there belongs also a finite, non-real $\omega$ (\S\,52). In the case of the
rational functions of $f(\omega)$, however, from the finiteness
for every finite imaginary $\omega$ we can only conclude that they are Laurent
polynomials in $f(\omega)$, allowing negative powers represented by $1:f(\omega)$,
because only to every finite
$f(\omega)$, with the exception of $f(\omega) = 0$, does there belong a finite
imaginary $\omega$.

But to the substitution

\weq{2}{\left(\omega, \frac{\omega-1}{\omega+1}\right)}

there corresponds the interchange

\[
\left(f(\omega), \frac{\sqrt{2}}{f(\omega)}\right), \qquad [\S\,34, (18)]
\]

and if, then, the functions $\Phi_{a,c,\partial}$ are so chosen that they
are permuted only among one another by the substitution (2) as well,
then the coefficients of the equation of which they are the roots will
depend rationally on
%CONT>>
